\chapter{関連研究}
\section{Human Sensing for Smart Cities}
本研究では，物理センサの課題を抱えるスマートシテ ィにおいて，市民を「人間センサ」として活用する方法
を提案している．人々がソーシャルメディアで共有する 認識を，確率的言語モデルを用いて抽出する手法を提示
している．ニューヨーク市の位置情報付きツイートを使 用し，電力管理，嵐への対応，公共安全といったスマー
トシティの取り組みに，人間センサが補完的，代替的， または補助的な情報源となり得ることを実証している．

\section{大規模災害時における被災者の感情を考慮した行 動促進ツイートの特徴分析}
本研究は，豪雨災害時における旧 Twitter のツイート
を対象に，交通施策に対する住民感情を分析した．自然
言語処理を用い，ツイート本文を形態素解析（Web 茶ま
め・MeCab）し，ツイートの感情極性値を定量的に算出
した．さらに，災害発生からの時期や交通体系の変化に 応じた 6
つのフェーズを設定し，各フェーズでの住民感
情の時系列変化を評価した．分析の結果，住民の感情は
交通ネットワークの状況変化に伴って時系列で変化する
ことが明らかになった．これにより SNS 上の情報が大
規模災害発生後の交通施策を講じる際の意思決定マネジ
メントにおいて有益な参照情報となることを示唆した．

\section{TODO: YouTubeライブ配信コメントを利用した研究何か}

\section{TODO: 他の参考文献}

\section{}